\begin{figure}[htbp]
\centering
\begin{tikzpicture}[
  >=Stealth,
  every node/.style={font=\small},
  entry/.style={inner sep=5pt},
  arrowlabel/.style={fill=white,inner sep=2pt,align=center}
]
  \draw[rounded corners,black!35] (-0.55,-1.95) rectangle (10.25,1.25);
  \node[anchor=west,fill=white,inner sep=3pt] at (-0.3,1.25)
    {Row $6$};
  \node[entry] (y) at (0.4,0) {$\mathsf y$};
  \node[entry,align=center] (p2) at (4.2,0.65)
    {$\mathsf p_2$\\[-1pt]\scriptsize detector for $b_2$};
  \node[entry,align=center] (p0) at (4.2,-0.85)
    {$\mathsf p_0$\\[-1pt]\scriptsize detectors for $b_{\jmath}$, $\jmath\ge4$};
  \node[draw,rounded corners,align=center,inner sep=6pt] (out) at (8.8,-0.85)
    {one of\\[2pt]
     $d_{85}:\ \mathsf c_2$\\[-1pt]
     $d_{157}:\ \mathsf a_1$};
  \draw[->,thick] (y) -- node[arrowlabel,above,sloped]
    {$d_{93}$ forced} (p2);
  \draw[->,dashed] (y) -- (p0);
  \node[arrowlabel] at (1.5,-1.25) {$d_{309}$ unavailable};
  \draw[->,thick] (p0) -- node[arrowlabel,above]
    {outgoing forced} (out);
\end{tikzpicture}
\caption{The final constraint in row $6$.  The known survivor $b_2$
forces the $d_{93}$ pairing, so its source $\mathsf y$ cannot also
hit $\mathsf p_0$ at length $309$.  The possible $\mathsf x$ incoming
at length $273$ lies in row $3$.  Thus $\mathsf p_0$ must support
one of the two indicated outgoing differentials; the argument does
not choose between them.  The diagram records logical constraints,
not chronological order: length $85$ precedes length $93$.}
\label{g:forcing-figure}
\end{figure}
